\begin{center}
    \Large \textbf{Abstract}
\end{center}

\vspace{0.8cm}

\addcontentsline{toc}{section}{Abstract}

Vehicle-to-Vehicle (V2V) communication is a critical enabler for cooperative driver 
assistance, yet established technologies such as DSRC and C-V2X presuppose infrastructure 
support that may be unavailable, sparse, or financially infeasible in resource-constrained 
deployment scenarios. This paper addresses the design and validation of a fully 
decentralized LoRa-based V2V mesh system capable of reliable multi-hop safety-message 
dissemination using only peer-to-peer forwarding and local decision-making, without 
dependence on roadside units, fixed gateways, or centralized coordination. While LoRa 
is not proposed as a replacement for purpose-built vehicular standards, its low cost
and accessibility make it a compelling experimental platform for studying decentralized 
dissemination under constrained sub-GHz channel conditions. The proposed system is 
organized as a three-layer architecture comprising an Application Layer for service-level 
logic, a Network Layer implementing managed flooding with duplicate suppression, 
hop-bounded propagation, and signal-quality-weighted relay timing, and a Link Layer 
providing channel-aware transmission through listen-before-send access control and 
selective acknowledgment-based reliability.

Vehicle-to-Vehicle (V2V) communication helps vehicles share safety information for
cooperative driver assistance. However, common options such as DSRC and C-V2X often
assume infrastructure support, which can be missing, limited, or too costly in
resource-constrained deployment scenarios. In this work, we design and validate a fully
decentralized LoRa-based V2V mesh system that can disseminate safety messages over
multiple hops using only peer-to-peer forwarding and local decision-making, without
roadside units, fixed gateways, or centralized coordination. We are not proposing LoRa as a
replacement for purpose-built vehicular standards; instead, we use it as a low-cost and
accessible experimental platform to study decentralized dissemination under constrained
sub-GHz channel conditions. The system follows a three-layer architecture: an Application
Layer for service-level logic, a Network Layer implementing managed flooding with
duplicate suppression, hop-bounded propagation, and signal-quality-weighted relay timing,
and a Link Layer providing channel-aware transmission through listen-before-send access
control and selective acknowledgment-based reliability.

We validate the design using two complementary methods: a deterministic host-side
simulation with virtual time progression and configurable channel impairments to verify
forwarding semantics under reproducible conditions, and real-world experiments with four
LoRa nodes demonstrating multi-hop routing, QoS-aware priority differentiation,
directional geographic propagation, and autonomous relay suppression. The results confirm
end-to-end message delivery across nodes beyond direct transmission range, correct priority
ordering with an observed delay ratio consistent with the intended fourfold scaling, a 33\%
reduction in relay events under directional mode, and distributed best-relay selection
without explicit inter-node negotiation. Overall, these findings show that a low-cost
LoRa-based platform can support prototyping and studying decentralized vehicular
networking principles in infrastructure-free environments.

\vspace{0.8cm}

\noindent\textbf{Keywords:} Vehicle-to-vehicle communication, LoRa mesh networks, cooperative safety, decentralized networking, intelligent transportation systems.